% ============================================================
% TEMPLATE: Projektdokumentation (Hochschule Düsseldorf)
% ------------------------------------------------------------
% Kompilieren:  pdflatex -> biber -> pdflatex -> pdflatex
% Benötigt:     lib.bib (Literatur), optional Logo-Datei
%
% So benutzt du das Template:
%   1. Block "METADATEN" ausfüllen (Titel, Autoren, Modul, ...)
%   2. Kapitel ab "INHALT" ersetzen
%   3. KI-Nutzung im Text mit \KIfn{...}{...}{...} ausweisen --
%      das KI-Verzeichnis (Tabelle am Ende) füllt sich automatisch.
% ============================================================
\documentclass[12pt,a4paper]{article}

% ============================================================
% PAKETE
% ============================================================
\usepackage[utf8]{inputenc}
\usepackage[T1]{fontenc}
\usepackage[ngerman]{babel}
\usepackage[left=3cm, right=3cm, top=2.5cm, bottom=2.5cm]{geometry}

\usepackage{graphicx}
\usepackage{float}

\usepackage{listings}
\usepackage{xcolor}

\usepackage{parskip}
\usepackage{enumitem}
\usepackage{booktabs}
\usepackage{longtable}
\usepackage{array}

\usepackage{hyperref}

% APA-Zitierweise (biblatex-apa) -- mit biber kompilieren!
\usepackage[style=apa, backend=biber]{biblatex}
\addbibresource{lib.bib}

% ============================================================
% METADATEN  <-- HIER ANPASSEN
% ============================================================
\newcommand{\hochschule}{Hochschule Düsseldorf}
\newcommand{\fachbereich}{Fachbereich Medien}
\newcommand{\adresse}{Münsterstraße 156\\40476 Düsseldorf}
\newcommand{\logodatei}{HSD (1).png}      % leer lassen oder Datei fehlt -> kein Logo

\newcommand{\studiengang}{B.\,Sc.\ Data Science, AI und Intelligente Systeme (DAISY)}
\newcommand{\semester}{Sommersemester 2026}
\newcommand{\fachsemester}{2}

\newcommand{\dokumenttyp}{Projektdokumentation}
\newcommand{\titel}{Titel der Arbeit}
\newcommand{\untertitel}{Untertitel der Arbeit}
\newcommand{\modul}{Modulbezeichnung}

\newcommand{\autoren}{%
    Vorname Nachname (Matr.-Nr: 0000000)\\
    Vorname Nachname (Matr.-Nr: 0000000)\\
    Vorname Nachname (Matr.-Nr: 0000000)%
}
\newcommand{\betreuer}{Herr Prof.\ Vorname Nachname}

\hypersetup{
    pdftitle={\titel},
    colorlinks=true,
    linkcolor=black,
    citecolor=black,
    urlcolor=blue
}

% ============================================================
% EINSTELLUNGEN
% ============================================================
\lstset{
    basicstyle=\ttfamily\small,
    breaklines=true,
    frame=single,
    columns=fullflexible,
    keepspaces=true
}
\lstdefinestyle{yaml}{
    basicstyle=\ttfamily\footnotesize,
    breaklines=true,
    frame=single,
    columns=fullflexible,
    keepspaces=true,
    backgroundcolor=\color{gray!5}
}
\setlength{\parskip}{6pt}

% ============================================================
% KI-VERZEICHNIS: automatische Sammlung der Prompts
% ------------------------------------------------------------
% \KIfn{Tool}{Prompt}{Verwendungszweck}
%     -> Fußnote an der Textstelle + Zeile im KI-Verzeichnis
% \KIfntext{Tool}{Prompt}{Verwendungszweck}
%     -> wie oben, aber als \footnotetext (z.B. nach Abbildungen/
%        Listen/Floats, wo \footnote nicht funktioniert)
% ============================================================
\makeatletter
\newcommand{\kirows}{}
\newcommand{\KIregister}[2]{%
    \g@addto@macro\kirows{#1 & „#2" \\}%
}
\newcommand{\KItext}[3]{%
    KI-Unterstützung (#1): „#2" Verwendungszweck: #3. %
    Siehe KI-Verzeichnis, Tabelle~\ref{tab:ki-verzeichnis}.%
}
\newcommand{\KIfn}[3]{%
    \KIregister{#1}{#2}%
    \footnote{\KItext{#1}{#2}{#3}}%
}
\newcommand{\KIfntext}[3]{%
    \KIregister{#1}{#2}%
    \footnotetext{\KItext{#1}{#2}{#3}}%
}
\makeatother

% ============================================================
% DOKUMENT
% ============================================================
\begin{document}

% ------------------------------------------------------------
% TITELSEITE
% ------------------------------------------------------------
\begin{titlepage}
    \begin{center}

        \IfFileExists{\logodatei}{\includegraphics[width=4cm]{\logodatei}\\[1cm]}{}

        {\large \hochschule}\\
        \fachbereich\\
        \adresse

        \vfill

        \textbf{Studiengang}\\[0.2cm]
        \studiengang\\
        \semester

        \vfill

        {\large \dokumenttyp}\\[0.4cm]
        {\huge \textbf{\titel}}\\[0.3cm]
        {\large \untertitel}\\[0.3cm]

        \vfill

        \textbf{\modul}\\[0.5cm]

        \vfill

        Vorgelegt von:\\[0.3cm]
        \autoren\\[0.4cm]

        Fachsemester: \fachsemester\\[0.2cm]
        Betreuer: \betreuer

    \end{center}
\end{titlepage}

% ------------------------------------------------------------
% VERZEICHNISSE
% ------------------------------------------------------------
\newpage
\pagenumbering{Roman}
\tableofcontents

\newpage
\begingroup
\makeatletter
\section*{Abbildungs- und Tabellenverzeichnis}
\@starttoc{lof}
\bigskip
\@starttoc{lot}
\makeatother
\endgroup

\newpage
\pagenumbering{arabic}

% ============================================================
% INHALT  <-- AB HIER ERSETZEN
% ============================================================
\section{Einleitung}

Fließtext mit Zitat \parencite{beispielquelle2026} bzw.\ \textcite{beispielquelle2026}.
Eine KI-Nutzung wird direkt an der Stelle ausgewiesen.\KIfn{ChatGPT}{Beispiel-Prompt im Wortlaut.}{Wofür die Ausgabe verwendet wurde}

\subsection{Unterabschnitt}

\subsubsection{Unter-Unterabschnitt}

% --- Beispiel: Abbildung -------------------------------------
% \begin{figure}[H]
%     \centering
%     \includegraphics[width=0.9\linewidth]{bild.png}
%     \caption[Kurztitel für Verzeichnis]{Langtitel \parencite{quelle}}
%     \label{fig:beispiel}
% \end{figure}
% \KIfntext{Gemini}{Prompt.}{Zweck}   % Fußnote nach einem Float

% --- Beispiel: Tabelle ---------------------------------------
\begin{table}[H]
    \centering
    \caption{Beispieltabelle}
    \label{tab:beispiel}
    \begin{tabular}{@{}p{3cm}p{4.8cm}p{5.7cm}@{}}
        \toprule
        \textbf{Kriterium} & \textbf{Option A} & \textbf{Option B} \\
        \midrule
        Kriterium 1 & Wert A1 & Wert B1 \\
        Kriterium 2 & Wert A2 & Wert B2 \\
        \bottomrule
    \end{tabular}
\end{table}

% --- Beispiel: Aufzählungen ----------------------------------
\begin{enumerate}
    \item \textbf{Punkt 1:} Beschreibung.
    \item \textbf{Punkt 2:} Beschreibung.
\end{enumerate}

\begin{description}[leftmargin=3cm, style=nextline]
    \item[Begriff A] Erklärung.
    \item[Begriff B] Erklärung.
\end{description}

% --- Beispiel: Code ------------------------------------------
\begin{lstlisting}
docker swarm init
\end{lstlisting}

\begin{lstlisting}[style=yaml]
# beispiel.yaml
key: value
list:
  - eintrag
\end{lstlisting}

\section{Hauptteil}

\section{Fazit}

% ============================================================
% KI-VERZEICHNIS (wird automatisch befüllt)
% ============================================================
\section{KI-Verzeichnis}

Sämtliche im Rahmen dieser Ausarbeitung verwendeten KI-Tools sind, analog zu den Literaturquellen, direkt an der jeweiligen Textstelle als Fußnote ausgewiesen -- mit vollständigem Prompt-Wortlaut sowie dem jeweiligen Verwendungszweck. Zur besseren Übersicht sind alle Prompts zusätzlich in Tabelle \ref{tab:ki-verzeichnis} zusammengefasst. Die inhaltliche Prüfung und Einordnung der KI-Ausgaben erfolgte in jedem Fall durch die Autor:innen.

\begin{longtable}{@{}p{2.5cm}p{11.5cm}@{}}
    \caption{Übersicht der verwendeten KI-Prompts}\label{tab:ki-verzeichnis}\\
    \toprule
    \textbf{Tool} & \textbf{Prompt} \\
    \midrule
    \endfirsthead
    \toprule
    \textbf{Tool} & \textbf{Prompt} \\
    \midrule
    \endhead
    \midrule
    \multicolumn{2}{r}{\textit{Fortsetzung auf der nächsten Seite}} \\
    \midrule
    \endfoot
    \bottomrule
    \endlastfoot
    \kirows
\end{longtable}

% ------------------------------------------------------------
% LITERATURVERZEICHNIS
% ------------------------------------------------------------
\newpage
\printbibliography

\end{document}
